\section{Analysis}\label{sec:analysis}

\subsection{Where AP is lost, and how Val400 relates to Test100}\label{sec:thresholds}

\paragraph{Per threshold.} \Cref{fig:thresholds} (left) plots Test100 AP at each OKS threshold, together with the \texttt{paper}-protocol ceiling of a perfect model (\cref{sec:eval}).
Every model sits on the ceiling up to $t=0.75$. They separate only above it: the released model's \APstrict{} is \res{rel.fidip.test.apstrict}.
On this test set, AP is decided by localisation precision, not by gross failures.

\paragraph{Val400 against Test100.} \Cref{fig:thresholds} (right) shows Val400 against Test100 AP. Across the \res{vt.n} runs with both, the Spearman rank correlation is \res{vt.spearman}, and Test100 exceeds Val400 by \res{vt.gap.mean}~AP on average.
Test100's looser boxes explain part of that gap (\cref{sec:released}). The ranking transfers reasonably even though Test100 was chosen for complex poses.

\begin{figure}[t]
  \centering
  \begin{minipage}[b]{0.63\linewidth}\centering
    \GenFigure[width=\linewidth]{oks_thresholds}{AP at OKS 0.50--0.95: released FiDIP, retrain, best config; paper-protocol ceiling band}
  \end{minipage}\hfill
  \begin{minipage}[b]{0.33\linewidth}\centering
    \GenFigure[width=\linewidth]{val_vs_test}{Val400 AP vs Test100 AP, one point per run that has both}
  \end{minipage}
  \caption{Left: Test100 AP at each OKS threshold, for the \texttt{paper} protocol (with the ceiling band of a perfect model under the id-0 artifact) and, in the second sub-panel, the \texttt{fixed} protocol. Right: Val400 against Test100 AP for every run evaluated on both sets.}
  \label{fig:thresholds}\label{fig:valtest}
\end{figure}

\paragraph{Evaluation choices.} \Cref{tab:sensitivity} (\cref{app:attempts}) holds the released model's predictions fixed and changes one evaluation choice at a time. Each choice moves AP by about as much as FiDIP's published gain over fine-tuning (0.9):
The matching protocol moves AP by \res{sens.id0.delta}, and moving test0's score rank alone swings the \texttt{paper} AP from \res{id0.rank.lo} to \res{id0.rank.hi}. Test100-clean moves it by \res{sens.clean.delta}, validate500's boxes give \res{sens.box.v500.ap} instead of \res{sens.box.v100.ap}, the misplaced \code{test11} costs \res{test11.delta}, and the placeholder area inflates Val400 to \res{sens.area.const.ap} from \res{rel.fidip.val.ap}.

\subsection{Per-keypoint accuracy and failure modes}\label{sec:perkp}\label{sec:failures}

\paragraph{Per keypoint.} \Cref{fig:perkp} (left) and \cref{tab:perkp} break the error down by joint. \textbf{Hips are by far the hardest joints}: their median error is more than twice that of any other group, and fewer than half fall within 5\% of the box scale.
This fits infant anatomy: hip centres are hidden under diapers and clothing, and legs are often flexed against the trunk. The paper likewise found the lower body hardest for adult-trained models~\citep{huang2021fidip}.
Heads and wrists are the most precise.

\paragraph{Failure modes.} \code{eval/taxonomy.py} assigns failure modes with unit-tested rules (\cref{tab:failures} in \cref{app:attempts} gives counts and exact criteria).
\textbf{Rotation dominates} (\cref{fig:perkp}, right): \res{fail.sideways.n} of the 100 infants lie sideways, with a mean OKS deficit of \res{fail.sideways.deficit}\% against \res{fail.upright.deficit}\% for upright infants. Left/right swaps are rare (\res{fail.lrswap.n} image, a back view), no prediction is a whole-body mirror, \res{fail.occl.n} images are heavily occluded and \res{fail.inverted.n} have an inverted torso.
The failures are mostly a matter of degree: sideways and partly hidden infants get somewhat imprecise joints, not catastrophic errors. That is consistent with AP being decided at strict thresholds.
\Cref{fig:qual} shows the worst Test100 cases.

\begin{figure}[t]
  \centering
  \begin{minipage}[b]{0.5\linewidth}\centering
    \GenFigure[width=\linewidth]{per_keypoint}{per-keypoint normalised error and PCK@0.05 (Test100)}
  \end{minipage}\hfill
  \begin{minipage}[b]{0.46\linewidth}\centering
    \GenFigure[width=\linewidth]{error_vs_orientation}{error vs torso orientation (Test100)}
  \end{minipage}
  \caption{Left: per-keypoint error on Test100, normalised by $\sqrt{wh}$ of the GT box (the OKS scale). Right: error by in-plane torso orientation; sideways infants carry most of the error.}
  \label{fig:perkp}
\end{figure}


\begin{figure}[t]
  \centering
  \GenFigure[width=\linewidth,height=0.2\textheight,keepaspectratio]{qual_failure}{overlay grid: worst Test100 cases}
  \caption{The lowest-OKS Test100 cases for the released FiDIP model, with GT and predictions overlaid (legend in the figure). The best cases are shown in \cref{fig:qualsucc}.}
  \label{fig:qual}
\end{figure}

\subsection{Inference demo}\label{sec:demo}

Part~6 of the course guide asks for two demos. The command-line demo \code{tools/demo.py} follows the guide's interface: \code{--cfg}, \code{--image}, \code{--model}, \code{--output}, and optionally \code{--conf-threshold} and \code{--output-size}; it also accepts \code{--video}.
\code{tools/streamlit_demo.py} lets a user upload an image, set the confidence threshold with a slider, and view keypoints with their confidences. Both were verified on the cluster, the Streamlit app through its AppTest harness and a live health check.

\Cref{fig:demo} (appendix) shows one practical lesson. The same image is processed with the GT box and with a detector box. Mean keypoint confidence drops from \res{demo.conf.gt} to \res{demo.conf.det}, so box quality matters as much as the pose model once GT boxes are gone.

